\section{Identification of the physical critical point}
\label{sec:criticality}

The parameter furnished by Theorem~\ref{traj:critical-trajectory} was
chosen by its scale evolution. We first identify that parameter by torus
duality, and then show that it is the threshold for plus magnetization.
The second step uses the uniform local comparison from
Proposition~\ref{cross:local-comparison}; self-duality alone will not be
used as a criterion for criticality.

Fix a sufficiently small $\theta>0$. Throughout this section the bond
couplings and their FK odds are
\begin{equation}
 v_e(h,\omega):=e^{2K_e(h,\omega)}-1
       =\sqrt2\,e^{h+\theta\xi_e(\omega)},
 \qquad h\in\R.
 \label{crit:odds}
\end{equation}
The original signs $\xi_e$ are independent and symmetric. Auxiliary
random choices used in the scale construction are integrated out whenever
we discuss the original spin or FK law.

For a finite graph $G=(V,E)$ and a partition $\pi$ of its boundary
vertices, write $\phi_{G,h,\omega}^{\pi}$ for the random-cluster measure
\begin{equation}
 \phi_{G,h,\omega}^{\pi}(\eta)
 =\frac{1}{\mathcal Z_{G,h,\omega}^{\pi}}
       2^{k_\pi(\eta)}\prod_{e\in E}v_e(h,\omega)^{\eta_e},
 \qquad \eta\in\{0,1\}^{E}.
 \label{crit:fk}
\end{equation}
Here $k_\pi(\eta)$ counts open clusters after identifying the vertices
in each part of $\pi$. The superscript $\mathrm w$ means that all boundary
vertices are identified. Unless explicitly stated otherwise, a path uses
ordinary lattice edges and does not use these identifications. At fixed
environment the measures in \eqref{crit:fk} are monotonic, satisfy FKG,
and increase with each edge odds and with boundary wiring. These are the
usual finite-volume FK properties; their Edwards--Sokal coupling to the
Ising model is given in~\cite{ES}.
We write $\phi_{h,\omega}^{\mathrm w}$ for the infinite-volume wired
measure, obtained by the monotone limits of finite wired boxes.

We use the following form of the decision-tree inequality. If $A$ is an
increasing event and a decision tree determines $A$ under a finite
monotonic measure $\phi$, then
\begin{equation}
 \phi(A)\bigl(1-\phi(A)\bigr)
 \le \sum_{e\in E}\delta_e\,
       \operatorname{Cov}_{\phi}(\1_A,\eta_e),
 \label{crit:osss}
\end{equation}
where $\delta_e$ is the probability that the tree queries $e$
\cite[Theorem~1.1]{DCRT}. The inequality also holds after averaging over
an independent choice of tree. For \eqref{crit:fk}, differentiation at
fixed environment gives the exact identity
\begin{equation}
 \frac{\dd}{\dd h}\phi_{G,h,\omega}^{\pi}(A)
 =\sum_{e\in E}
       \operatorname{Cov}_{\phi_{G,h,\omega}^{\pi}}(\1_A,\eta_e).
 \label{crit:derivative}
\end{equation}
Indeed, varying $h$ multiplies the unnormalized weight by
$\exp(h\sum_e\eta_e)$, apart from its $h$-independent factors. Both
\eqref{crit:osss} and \eqref{crit:derivative} will be applied before
averaging over environments.

\subsection{Torus duality fixes the tuned parameter}

Let $\mathbb T_N=(\Z/N\Z)^2$, with $N\ge3$. A spin holonomy
$\alpha=(\alpha_1,\alpha_2)\in\{0,1\}^2$ prescribes the signs acquired
by a spin after one turn around the two coordinate cycles. Equivalently,
choose edge signs $\tau_e^\alpha$ whose product around every elementary
face is $1$ and whose products around the two coordinate cycles are
$(-1)^{\alpha_1}$ and $(-1)^{\alpha_2}$. Define
\[
 Z_\alpha(K)=\sum_{\sigma\in\{-1,1\}^{\mathbb T_N}}
       \exp\left(\sum_{e=\{x,y\}}K_e\tau_e^\alpha\sigma_x\sigma_y\right).
\]
Changing the representative $\tau^\alpha$ by a vertex gauge
transformation does not change this sum. Write $Z(K)$ for the vector of
these four partition functions, ordered as $00,10,01,11$, and put
$\widehat Z(K)=Z(K)/\|Z(K)\|_2$.

\begin{lemma}[The four torus partition functions]
\label{crit:torus-identities}
For positive couplings on $\mathbb T_N$, $N\ge3$, let the dual couplings satisfy
$e^{-2K_{e^*}^*}=\tanh K_e$. Identify the dual square torus with the
original one by a half-lattice translation. Then
\begin{equation}
 Z(K)=a(K)\mathsf H Z(K^*),\qquad
 a(K)=\prod_e\sqrt{\sinh(2K_e)},\qquad
 \mathsf H=\frac12
 \begin{pmatrix}
 1&1&1&1\\
 1&1&-1&-1\\
 1&-1&1&-1\\
 1&-1&-1&1
 \end{pmatrix}.
 \label{crit:torus-duality}
\end{equation}
In particular $\mathsf H$ is orthogonal and fixes the pure critical
vector. For every holonomy $\alpha$,
\begin{equation}
 \frac{Z_\alpha(K)}{Z_{00}(K)}
 =\phi_{\mathbb T_N,K}
 \left(\alpha([\gamma])=0
       \text{ for every open cycle }\gamma\right),
 \label{crit:torus-ratio}
\end{equation}
where $[\gamma]\in(\Z/2\Z)^2$ is the cycle's winding class and
$\phi_{\mathbb T_N,K}$ is the untwisted torus FK law.
\end{lemma}

\begin{proof}
We record the finite-volume form of Kramers--Wannier duality
\cite{KW}, including its homology convention; the four-sector matrix
also appears in \cite[Equations~(4)--(5)]{BS}. For
$b\in(\Z/2\Z)^2$, let
\[
 W_b(K)=\sum_{\substack{A\subset E(\mathbb T_N)\text{ even}\\{}[A]=b}}
                    \prod_{e\in A}\tanh K_e.
\]
The high-temperature expansion gives
\[
 Z_\alpha(K)=2^{N^2}\prod_e\cosh K_e
                  \sum_b(-1)^{\alpha\cdot b}W_b(K).
\]
The low-temperature expansion of the dual spin system with holonomy
$\beta$ has domain walls of winding class
$(\beta_2,\beta_1)$. Each admissible wall configuration comes from two
spin configurations, so
\[
 Z_\beta(K^*)=2e^{\sum_eK_{e^*}^*}
                   W_{(\beta_2,\beta_1)}(K).
\]
Substitution gives the sign
$(-1)^{\alpha_1\beta_2+\alpha_2\beta_1}$ in
\eqref{crit:torus-duality}. Its prefactor follows from
$|E(\mathbb T_N)|=2N^2$ and
$e^{-K_{e^*}^*}\cosh K_e=\sqrt{\sinh(2K_e)/2}$.
At $K_e=K_0$, the dual couplings equal $K_0$ and $a(K)=1$; hence
$\mathsf HZ(K_0)=Z(K_0)$. The displayed matrix satisfies
$\mathsf H^2=I$.

For the ratio, expand each edge factor as
\[
 e^{K_e\tau_e^\alpha\sigma_x\sigma_y}
 =e^{-K_e}\left(1+v_e\,
           \1_{\{\sigma_x=\tau_e^\alpha\sigma_y\}}\right).
\]
On a chosen open subgraph the spin constraints are consistent exactly
when every open cycle has zero evaluation under $\alpha$. In that case
there are $2^{k(A)}$ compatible spin assignments. Dividing the resulting
sum by its untwisted version proves \eqref{crit:torus-ratio}.
\end{proof}

\begin{proposition}[Identification of the tuned parameter]
\label{crit:self-dual-tuning}
The value $h_*$ furnished by
Theorem~\ref{traj:critical-trajectory} is zero.
\end{proposition}

\begin{proof}
Choose a fixed large integer $M$ and torus periods $N_k=Ms_k$, so that
every complete support and environment-dependence neighborhood of
Proposition~\ref{map:step}, through scale $k$, embeds in the torus.
The primitive environment variables in each such neighborhood have
their plane distribution, so each resulting potential has its plane
law when expressed in local spin coordinates. Apply the
scale construction with iid torus disorder, simultaneously in the four
spin holonomies. Let $\zeta$ denote the independent algorithmic choices
used in this run. In every contractible neighborhood a flat twist has a
spin trivialization; changes of trivialization on overlaps are constant
spin flips. Every interaction is even, and the auxiliary variables are
unchanged by spin flip. The local identities therefore hold in each
holonomy, with the same scalar factors removed by centering. In
particular, for a common positive factor $c_k(\omega,\zeta)$ the resulting
partition functions satisfy
\[
 e^{-S_k}Z_\alpha(K_0)
 \le c_k(\omega,\zeta)^{-1}Z_\alpha(K(h_*,\omega))
 \le e^{S_k}Z_\alpha(K_0),
\]
where $S_k=S_k(\omega,\zeta)$ can be the largest, over the four holonomies, of the sums of
norms of the remaining potentials. There are only $M^2$ cells at the
last scale. The trajectory bound $X_k\to0$ gives
$\E S_k\le C M^2 X_k\to0$. Consequently
\begin{equation}
 \bigl\|\widehat Z(K(h_*,\omega))-
                 \widehat Z(K_0)\bigr\|_2\longrightarrow0
 \quad\text{in environmental probability}.
 \label{crit:torus-terminal}
\end{equation}
The partition vector itself does not depend on the algorithmic choices,
so this conclusion also holds after retaining only the original bond
environment.

The dual odds satisfy $v_{e^*}^*=2/v_e$. Thus duality sends
$(h,\xi_e)$ to $(-h,-\xi_e)$, with the signs transported along the
crossing-edge bijection. This transformation preserves their iid law.
By Lemma~\ref{crit:torus-identities}, the normalized duality transform is
the orthogonal matrix $\mathsf H$, and its pure target is fixed.
Therefore \eqref{crit:torus-terminal} holds with $-h_*$ as well.

Take $\alpha=10$, and let $A_\alpha$ be the increasing event that some
open cycle has odd evaluation under $\alpha$. Its pure probability
$q_N^0$ satisfies
\begin{equation}
 c\le q_N^0\le1-c
 \label{crit:pure-torus-nondegeneracy}
\end{equation}
uniformly for large $N$. For the lower bound, a finite chain of pure FK
rectangle crossings in a torus band constructs an open cycle winding
once horizontally. For the upper bound, the corresponding construction
of a closed-dual cycle winding once vertically prevents such an open
cycle. The constructions use finitely many embedded rectangles of fixed
aspect ratios and the uniform pure FK bounds recalled in
\cite[Section~2, ``Pure buffered RSW consequences'']{S}.
They can be imposed successively with arbitrary induced boundary wiring;
for a closed-dual crossing use the worst, fully wired primal boundary.

All ratios in \eqref{crit:torus-ratio} belong to $[0,1]$, so the
untwisted coordinate of $\widehat Z$ is at least $1/2$. Hence
\eqref{crit:torus-terminal} implies, on the same original signs,
\[
 q_{N_k,\omega}(h_*)-q_{N_k}^0\longrightarrow0,
 \qquad
 q_{N_k,\omega}(-h_*)-q_{N_k}^0\longrightarrow0
\]
in probability, where $q_{N,\omega}(h)=
\phi_{\mathbb T_N,h,\omega}(A_\alpha)$. The two convergence statements
hold simultaneously by a union bound. Querying every edge in
\eqref{crit:osss} and using \eqref{crit:derivative} gives
$q'_{N,\omega}(h)\ge q_{N,\omega}(h)(1-q_{N,\omega}(h))$.
If $h_*\ne0$, monotonicity and
\eqref{crit:pure-torus-nondegeneracy} keep this product bounded below on
the interval between $-h_*$ and $h_*$, with probability tending to one.
Integration forces an endpoint difference bounded below by a positive
constant, whereas both endpoints approach the same $q_{N_k}^0$.
This contradiction proves $h_*=0$.
\end{proof}

\subsection{From local barriers to sharpness}

We now use the local crossing comparison at $h=0$. Put
$B_r(x)=\{y\in\Z^2:|y-x|_\infty\le r\}$ and $B_r=B_r(0)$.
All annuli below have deterministic lattice boundaries; changing them by
a bounded incident layer changes none of the estimates.

\begin{lemma}[Local FK barriers]
\label{crit:ring-barrier}
There are fixed $a>1$ and $c_0>0$ with the following property. Let
$\mathcal A(x,r)$ be the lattice annulus
$B_{ar}(x)\setminus B_r(x)$, and let $\mathcal R(x,r)$ be its ordinary
radial crossing event. Call its bond environment good if
\[
 \phi_{\mathcal A(x,r),0,\omega}^{\mathrm w}
       (\mathcal R(x,r))\le1-c_0,
\]
where both boundary components are wired together. The probability that the
environment is not good tends to zero as $r\to\infty$, uniformly in
$x$. Goodness is a function only of the bonds of that annulus and its
incident layer. Every good annulus has the same crossing upper bound
for all $h\le0$ and all less wired boundary conditions. In particular,
\begin{equation}
 \E\phi_{0,\omega}^{\mathrm w}(0\leftrightarrow\infty)=0.
 \label{crit:no-percolation-at-tuning}
\end{equation}
\end{lemma}

\begin{proof}
Choose two separated buffered subbands of an annulus of fixed sufficiently
large ratio $a$. The pure probability of simultaneous plus and minus spin
circuits, one in each subband, is bounded below by a constant $2c_0>0$,
uniformly over pins outside the buffers; this is precisely
\cite[Section~2, Proposition ``Pure buffered RSW consequences'']{S}.
Proposition~\ref{cross:local-comparison} transfers this lower bound to
the random spin law, with probability tending to one over the local
environment and uniformly over the boundary pins. Apply it to
Edwards--Sokal spins with the wired boundary cluster pinned plus.
An ordinary open FK path has a constant spin. Such a path cannot cross
both separating circuits of opposite colors. The common triangulation
used for spin circuits causes no exception: a nearest-neighbor lattice
path can cross a circuit only through a vertex of that circuit. The
ordinary radial FK crossing probability is thus at most $1-c_0$.

Defining goodness by this wired crossing probability makes it a function
of the original local bonds alone. Its failure probability tends to zero
after integrating out all algorithmic choices. FK monotonicity gives
the assertion for $h\le0$ and for weaker wiring.

For completeness, this also rules out percolation at $h=0$. Pack
concentric annuli with a fixed large geometric separation, so that their
bond sets, including incident layers, are disjoint. Conditional on all
edges off one annulus, its induced FK boundary partition is dominated by
full wiring. The probability of an arm crossing $m$ such annuli is
therefore bounded, at fixed environment, by the product of their fully
wired crossing bounds. Choose the first radius large enough that every
bad-annulus probability is at most $1/2$. Independence of the local bond
sets bounds the environmental mean of this product by
$(1-c_0/2)^m$. Letting $m\to\infty$ in the wired infinite-volume limit
proves \eqref{crit:no-percolation-at-tuning}.
\end{proof}

The preceding lemma gives no rate for the probability of a bad annulus.
The next proof obtains the rate needed for sharpness by using many
independent annuli, before applying the decision-tree inequality.

\begin{lemma}[Subcritical arms]
\label{crit:subcritical-arms}
For every $b>0$ and $A>0$ there is $C_{b,A}<\infty$ such that
\begin{equation}
 \E\phi_{B_{2N}(x),-b,\omega}^{\mathrm w}
 \bigl(x\leftrightarrow\partial B_N(x)
                \text{ inside }B_N(x)\bigr)
       \le C_{b,A}N^{-A}.
 \label{crit:annealed-arm}
\end{equation}
The same bound holds for the ordinary path event under any finite or
infinite-volume FK law whose domain contains $B_{2N}(x)$, with arbitrary
boundary condition outside that box.
\end{lemma}

\begin{proof}
We work at $x=0$; the estimates are uniform over translates. For every
vertex $z\in B_{N+1}$, pack $m_N\ge c\log N$ separated annuli centered
at $z$, with inner radii between $N^{1/4}$ and a sufficiently small
multiple of $N^{1/2}$. Their full annular graphs lie in
$B_{N^{1/2}}(z)$, and their bond neighborhoods are disjoint. Let
$\eta_N\to0$ be an upper bound for their bad-environment probabilities,
uniformly over all these radii and centers, provided by
Lemma~\ref{crit:ring-barrier}. Independence for the annuli around one
fixed $z$ gives
\[
 \Prob\bigl(\text{at least }m_N/2\text{ annuli at }z
                      \text{ are bad}\bigr)
       \le 2^{m_N}\eta_N^{m_N/2}.
\]
Let $\mathcal G_N$ be the event that fewer than half are bad at every
$z\in B_{N+1}$. A union bound, without any independence assertion
between distinct centers, gives
\begin{equation}
 \Prob(\mathcal G_N^c)
 \le C N^2 2^{m_N}\eta_N^{m_N/2}
       =o(N^{-A})\qquad\text{for every }A>0.
 \label{crit:environment-exception}
\end{equation}
Indeed $m_N$ is bounded above and below by constant multiples of
$\log N$, while $\log(1/\eta_N)\to\infty$.

On $\mathcal G_N$, multiplication of the good-annulus bounds gives,
for some fixed $\alpha>0$,
\begin{equation}
 \phi_{B_{2N},h,\omega}^{\mathrm w}
 \bigl(z\leftrightarrow\partial B_{N^{1/2}}(z)
                   \text{ inside }B_{N^{1/2}}(z)\bigr)
       \le C N^{-\alpha},
 \qquad z\in B_{N+1},\quad h\le0.
 \label{crit:uniform-short-arm}
\end{equation}
For large $N$ all these small boxes lie within $B_{2N}$. Induced
boundary conditions are harmless because the annular bounds already
use full wiring. The event $\mathcal G_N$ works simultaneously for
every $h\le0$ by monotonicity.

Fix such an environment and write
\[
 q_\omega(h)=\phi_{B_{2N},h,\omega}^{\mathrm w}
     (0\leftrightarrow\partial B_N\text{ inside }B_N).
\]
Choose an integer $\ell$ uniformly between $N/4$ and $N/2$, and
explore all ordinary open clusters in $B_N$ that meet the layer
$\partial B_\ell$. This determines the event defining $q_\omega$:
every path from $0$ to $\partial B_N$ meets that layer. The exploration
uses only ordinary edges in $B_N$, although its law is the marginal
of the wired measure on $B_{2N}$. This is the layer exploration of
\cite[Lemma~3.2 and its proof]{DCRT}.

An edge is queried only if an endpoint lies on the chosen layer or
has an ordinary open path to it. For a fixed endpoint $z$, the
probability, over $\ell$, that the layer is within distance
$N^{1/2}$ of $z$ is at most $C N^{-1/2}$. Otherwise such a path forces
an arm of length $N^{1/2}$ at $z$. Thus
\eqref{crit:uniform-short-arm} bounds the maximum averaged revealment
by
\[
 \max_e\delta_e\le C\bigl(N^{-1/2}+N^{-\alpha}\bigr)
                   \le C N^{-d},\qquad
 d=\min\{1/2,\alpha\}>0.
\]
Edges outside $B_N$ have revealment zero. Their covariances are
nonnegative, so retaining them in the derivative sum only strengthens
the bound obtained from \eqref{crit:osss}. Equations
\eqref{crit:osss}--\eqref{crit:derivative} yield, for all $h\le0$,
\begin{equation}
 q'_\omega(h)\ge cN^d q_\omega(h)(1-q_\omega(h)).
 \label{crit:sharp-differential}
\end{equation}
Moreover \eqref{crit:uniform-short-arm} at $z=0$ gives
$q_\omega(0)\le C N^{-\alpha}$. For sufficiently large $N$, integrating
the derivative of $\log(q_\omega/(1-q_\omega))$ over $[-b,0]$ gives
\[
 \frac{q_\omega(-b)}{1-q_\omega(-b)}
 \le \frac{q_\omega(0)}{1-q_\omega(0)}e^{-cbN^d}
 \le e^{-cbN^d}.
\]
Combining this estimate with \eqref{crit:environment-exception} proves
\eqref{crit:annealed-arm}. In a larger graph, condition outside
$B_{2N}(x)$ and dominate the induced boundary partition by full
wiring. This proves the final assertion, including infinite-volume
limits.
\end{proof}

\subsection{The magnetization threshold}

The preceding results now identify the physical temperature. Put
$m_\theta(h)=\E\phi_{h,\omega}^{\mathrm w}(0\leftrightarrow\infty)$.
Edwards--Sokal identifies this with the environmental mean of the
infinite-volume plus magnetization. Under the plus boundary convention of
Section~\ref{sec:introduction}, the finite FK graph contains all bonds meeting the spin box,
and all exterior endpoints are wired and assigned spin plus. Shifting
the wired boundary by one lattice layer sandwiches this law between
the usual wired-box laws, so their infinite-volume limits agree.

\begin{theorem}[Physical critical point]
\label{crit:physical}
For every sufficiently small fixed $\theta>0$,
\[
 m_\theta(h)=0\quad(h\le0),\qquad
 m_\theta(h)>0\quad(h>0).
\]
Consequently, for every sufficiently small $\epsilon>0$, the critical
inverse temperature in Theorem~\ref{thm:main} is the unique positive
solution of
\begin{equation}
 \bigl(e^{2\beta_c(\epsilon)(1+\epsilon)}-1\bigr)
 \bigl(e^{2\beta_c(\epsilon)(1-\epsilon)}-1\bigr)=2.
 \label{crit:physical-equation}
\end{equation}
The plus magnetization is zero at this temperature.
\end{theorem}

\begin{proof}
Lemma~\ref{crit:ring-barrier} and monotonicity give $m_\theta(h)=0$
for $h\le0$. Fix $h=b>0$. Planar FK duality replaces the edge odds
by $2/v_e$, so the dual environment has the law of the model at $-b$.
In finite domains the dual boundary condition is dominated by full
wiring. Taking limits, the dual of the infinite-volume wired primal
law therefore satisfies the local annealed arm bounds of
Lemma~\ref{crit:subcritical-arms}.

A closed-dual circuit surrounding $B_M$ whose farthest distance from
the origin lies in $[R,2R)$ has diameter at least $cR$. It consequently
contains an ordinary dual arm of length $c'R$ from some dual vertex
in $B_{2R}$, with absolute positive constants $c,c'$. There are
$O(R^2)$ possible vertices. Using \eqref{crit:annealed-arm} with any
$A>2$ and summing over dyadic $R\ge c''M$ gives
\begin{equation}
 \E\phi_{b,\omega}^{\mathrm w}
 \bigl(\text{a closed-dual circuit surrounds }B_M\bigr)
 \le C_{b,A}\sum_{j\ge0}(2^jM)^{2-A}
 \longrightarrow0.
 \label{crit:dual-circuits}
\end{equation}
Here the arm bound is applied in a box of radius comparable with
$R$, with a further collar of comparable width; its validity for
arbitrary exterior boundary conditions was part of
Lemma~\ref{crit:subcritical-arms}.

If every open cluster meeting $B_M$ were finite, their union would
be finite. It contains $B_M$, and its lattice edge boundary consists
of closed edges. Connecting the union through the box, if necessary,
does not change its exterior boundary. Its outer dual boundary
therefore contains a closed-dual circuit surrounding $B_M$.
The complement of the event in \eqref{crit:dual-circuits} thus forces
some vertex of $B_M$ to belong to an infinite open cluster. For large
$M$ this has positive joint probability. The joint environment--wired
FK law is translation invariant, whence
\[
 0<\E\phi_{b,\omega}^{\mathrm w}(B_M\leftrightarrow\infty)
 \le |B_M|\,m_\theta(b).
\]
This proves positivity for every $b>0$.

Finally, the left side of \eqref{crit:physical-equation}, viewed as
a function of $\beta\ge0$, is continuous and strictly increasing
from zero to infinity. Let $\beta_*$ be its unique solution, and set
\[
 v_\pm=e^{2\beta_*(1\pm\epsilon)}-1,
 \qquad
 \theta_\epsilon=\frac12\log\frac{v_+}{v_-}.
\]
Then $v_\pm=\sqrt2e^{\pm\theta_\epsilon}$. Comparison with the
constant couplings gives
$K_0/(1+\epsilon)\le\beta_*\le K_0/(1-\epsilon)$, so
$\theta_\epsilon\to0$ as $\epsilon\to0$. The model at $\beta_*$
therefore falls within the result just proved and has zero plus
magnetization. The same holds below $\beta_*$ by monotonicity.
For $\beta>\beta_*$ both log odds increase, and
\[
 b:=\min_{s\in\{-1,1\}}
       \log\frac{e^{2\beta(1+s\epsilon)}-1}
                    {e^{2\beta_*(1+s\epsilon)}-1}>0.
\]
The bonds at $\beta$ dominate those of
\eqref{crit:odds} with parameters $(b,\theta_\epsilon)$ in the
same sign environment. Their plus magnetization is consequently
positive. This proves $\beta_*=\beta_c(\epsilon)$ with the
normalization and plus boundary convention of Theorem~\ref{thm:main}.
\end{proof}
